% !TEX program = xelatex
% !TEX root = ../../TFM.tex
\documentclass[../../TFM.tex]{subfiles}

\begin{document}

	% =====================================================
	% ANEXOS
	% =====================================================
	\appendix
	\chapter{Anexos}
	\label{chap:anexos}
	\thesisstartchap{anexos}

	Este capítulo reúne materiales complementarios de apoyo empírico, gráfico y
	descriptivo. Su propósito es documentar las fuentes y la muestra, sintetizar
	el diseño analítico y presentar los diagnósticos y sensibilidades que respaldan
	la lectura de las estimaciones principales.

	\section{Fuentes de datos e indicadores}
	\label{sec:fuentes}

	La \autoref{tab:fuentes_datos} documenta la definición activa de cada variable, la fuente efectivamente incorporada al panel, el indicador de origen y las transformaciones aplicadas. Los períodos corresponden a los años con información disponible en los insumos conservados en \texttt{data/raw}; la última columna informa la cobertura observada en la grilla maestra de 55 países y 1.430 observaciones país--año para 1996--2021. Las fuentes mantenidas únicamente para contrastar definiciones se identifican en la nota y no se combinan con las series activas.

	\renewcommand{\arraystretch}{1.10}
	\setlength{\tabcolsep}{3pt}

	\begin{footnotesize}
		\begin{longtable}{
				L{2.8cm}
				L{3.65cm}
				L{3.85cm}
				L{1.35cm}
				>{\RaggedLeft\arraybackslash}p{1.65cm}
			}

			\caption[Fuentes e indicadores del panel empírico]{Fuentes, indicadores y transformaciones utilizados en el panel empírico.}

			\label{tab:fuentes_datos} \\
			\hline
			Variable & Definición activa & Fuente efectiva e indicador de origen & Período & Cobertura \\
			\hline
			\endfirsthead

			\hline
			Variable & Definición activa & Fuente efectiva e indicador de origen & Período & Cobertura \\
			\hline
			\endhead

			$DRES$
			& Participación de las exportaciones de recursos no renovables del subsuelo en las exportaciones de mercancías. Incluye SITC Rev.\ 2: sección 3, divisiones 27, 28 y 68, partidas 6672--6673 y 9710; excluye los códigos no extractivos 2711, 2786, 2820, 2881, 2882, 3510, 6671 y 6674. El denominador incluye \texttt{XXXX} y excluye servicios.
			& Harvard Growth Lab, \textit{Atlas of Economic Complexity}, comercio bilateral SITC Rev.\ 2 a cuatro dígitos; construcción propia.
			& 1990--2022; base 1990--1995
			& 55 países \\

			$ECI$
			& Índice de Complejidad Económica calculado por el Atlas bajo la clasificación HS92.
			& Harvard Growth Lab, \textit{Atlas of Economic Complexity}, API GraphQL, campo \texttt{countryYear.eci}.
			& 1995--2022
			& 1.430 (100,00\,\%) \\

			$HHI$
			& Suma de los cuadrados de las participaciones de productos SITC Rev.\ 2 a cuatro dígitos; excluye servicios y el residual \texttt{XXXX}.
			& Harvard Growth Lab, \textit{Atlas of Economic Complexity}; cálculo propio a partir del comercio por producto.
			& 1990--2022
			& 1.430 (100,00\,\%) \\

			$DIVX$
			& Diversificación exportadora, definida como $1-HHI$.
			& Harvard Growth Lab, \textit{Atlas of Economic Complexity}; cálculo propio a partir de $HHI$.
			& 1990--2022
			& 1.430 (100,00\,\%) \\

			$RENTS$
			& Suma de las rentas de petróleo, gas natural, carbón y minerales como porcentaje del PIB, únicamente cuando están observados los cuatro componentes.
			& Banco Mundial, WDI: \texttt{NY.GDP.PETR.RT.ZS}, \texttt{NY.GDP.NGAS.RT.ZS}, \texttt{NY.GDP.COAL.RT.ZS} y \texttt{NY.GDP.MINR.RT.ZS}.
			& 1980--2021
			& 1.352 (94,55\,\%) \\

			$RENTS^{OIL+GAS}$
			& Suma de las rentas de petróleo y gas natural como porcentaje del PIB cuando ambos componentes están observados.
			& Banco Mundial, WDI: \texttt{NY.GDP.PETR.RT.ZS} y \texttt{NY.GDP.NGAS.RT.ZS}; cálculo propio.
			& 1980--2021
			& 1.365 (95,45\,\%) \\

			$RENTS^{MINING}$
			& Suma de las rentas de carbón y minerales como porcentaje del PIB cuando ambos componentes están observados.
			& Banco Mundial, WDI: \texttt{NY.GDP.COAL.RT.ZS} y \texttt{NY.GDP.MINR.RT.ZS}; cálculo propio.
			& 1980--2021
			& 1.384 (96,78\,\%) \\

			$INST$
			& Promedio simple de Control de la Corrupción, Estado de Derecho y Efectividad Gubernamental, solo cuando los tres componentes están observados.
			& Banco Mundial, WGI: \texttt{GOV\_WGI\_CC.EST}, \texttt{GOV\_WGI\_RL.EST} y \texttt{GOV\_WGI\_GE.EST}.
			& 1996--2024
			& 1.257 (87,90\,\%) \\

			$RENTS\times INST$
			& Producto directo de $RENTS$ e $INST$ cuando ambas variables están observadas.
			& Elaboración propia a partir de las dos series activas en el panel maestro.
			& 1996--2021
			& 1.197 (83,71\,\%) \\

			$OILPC$
			& Renta petrolera real por habitante: renta petrolera sobre PIB multiplicada por el PIB real y dividida por la población.
			& Banco Mundial, WDI: \texttt{NY.GDP.PETR.RT.ZS}, \texttt{NY.GDP.MKTP.KD} y \texttt{SP.POP.TOTL}; cálculo propio.
			& 1980--2021
			& 1.397 (97,69\,\%) \\

			$GASPC$
			& Renta de gas natural real por habitante, construida de forma análoga a $OILPC$.
			& Banco Mundial, WDI: \texttt{NY.GDP.NGAS.RT.ZS}, \texttt{NY.GDP.MKTP.KD} y \texttt{SP.POP.TOTL}; cálculo propio.
			& 1980--2021
			& 1.365 (95,45\,\%) \\

			$COALPC$
			& Renta del carbón real por habitante, construida de forma análoga a $OILPC$.
			& Banco Mundial, WDI: \texttt{NY.GDP.COAL.RT.ZS}, \texttt{NY.GDP.MKTP.KD} y \texttt{SP.POP.TOTL}; cálculo propio.
			& 1980--2021
			& 1.384 (96,78\,\%) \\

			$PEXP$
			& Participación de las secciones SITC Rev.\ 2 0, 1, 2 y 4 y de la división 68 en las exportaciones de mercancías.
			& Harvard Growth Lab, \textit{Atlas of Economic Complexity}, comercio SITC Rev.\ 2 a cuatro dígitos; cálculo propio.
			& 1990--2022
			& 1.430 (100,00\,\%) \\

			$FEXP$
			& Participación de la sección SITC Rev.\ 2 3 en las exportaciones de mercancías.
			& Harvard Growth Lab, \textit{Atlas of Economic Complexity}, comercio SITC Rev.\ 2 a cuatro dígitos; cálculo propio.
			& 1990--2022
			& 1.430 (100,00\,\%) \\

			$VOL$
			& Desviación estándar muestral móvil de cinco años de $100\Delta\log(CEMPI)$; cada valor requiere seis niveles anuales consecutivos.
			& FMI, \textit{Commodity Terms of Trade}: \texttt{CEMPI\_CTOTNX\_TT}, ponderaciones fijas \texttt{H\_FW\_IX}, frecuencia anual.
			& 1962--2022
			& 1.403 (98,11\,\%) \\

			$RER$
			& Logaritmo natural del nivel de precios del PIB de producción; un aumento representa una apreciación real.
			& Penn World Table 11.0, variable \texttt{pl\_gdpo}.
			& 1950--2023
			& 1.352 (94,55\,\%) \\

			$HUMCAP$
			& Exponencial de la función por tramos de retornos educativos de PWT aplicada a los años promedio de escolaridad de la población de 25 años o más.
			& PNUD, \textit{Human Development Report 2025}, \texttt{UNDP\_HDR\_MYS}; transformación propia.
			& 1990--2023
			& 1.403 (98,11\,\%) \\

			$INNOV$
			& Logaritmo natural de uno más el número de artículos científicos y técnicos por millón de habitantes.
			& Banco Mundial, WDI: \texttt{IP.JRN.ARTC.SC} y \texttt{SP.POP.TOTL}; cálculo propio.
			& 1996--2022
			& 1.430 (100,00\,\%) \\

			$NET$
			& Personas que utilizan internet como porcentaje de la población.
			& Banco Mundial, WDI: \texttt{IT.NET.USER.ZS}.
			& 1990--2022
			& 1.379 (96,43\,\%) \\

			$LOG\_GDPPC$
			& Logaritmo natural del PIB per cápita a PPA en precios constantes.
			& Banco Mundial, WDI: \texttt{NY.GDP.PCAP.PP.KD}.
			& 1990--2022
			& 1.378 (96,36\,\%) \\

			$GOVCONS$
			& Gasto de consumo final del gobierno general como porcentaje del PIB.
			& División de Estadística de las Naciones Unidas, \textit{National Accounts Main Aggregates}, serie 16, componente 3.
			& 1970--2024
			& 1.430 (100,00\,\%) \\

			$FIN$
			& Crédito doméstico al sector privado otorgado por bancos como porcentaje del PIB; proxy de profundidad del crédito bancario.
			& Banco Mundial, WDI: \texttt{FD.AST.PRVT.GD.ZS}.
			& 1980--2025
			& 1.328 (92,87\,\%) \\

			\hline
			\multicolumn{5}{p{14.45cm}}{\textit{Nota: La cobertura se expresa como observaciones disponibles y porcentaje de la grilla maestra 1996--2021, salvo $DRES$, que se calcula como promedio 1990--1995 y se utiliza exclusivamente para seleccionar los 55 países mediante el umbral del 20\,\%. No se imputan faltantes ni se rellenan vacíos mediante fuentes alternativas. La descarga WDI asociada a $DRES$ se conserva únicamente para validación agregada y no interviene en su cálculo. Para $GOVCONS$, WDI \texttt{NE.CON.GOVT.ZS} se conserva solo como contraste frente a la serie activa de Naciones Unidas. Para $FIN$, WDI \texttt{FS.AST.PRVT.GD.ZS} se conserva solo como contraste frente a la serie bancaria activa. Para $HUMCAP$, \texttt{hc} de PWT 11.0 se conserva solo como contraste frente al índice activo construido con escolaridad del PNUD. $HHI$ se utiliza en las ecuaciones con $ECI$ y se excluye cuando $DIVX=1-HHI$ es la variable dependiente.}} \\
		\end{longtable}
	\end{footnotesize}
	\thesissource{Elaboración propia a partir de los archivos, metadatos y scripts reproducibles del proyecto.}

	\section{Países seleccionados mediante DRES}
	\label{sec:muestra-dres}

	La \autoref{tab:muestra_dres20} identifica las 55 economías que integran la
	grilla maestra. La selección exige que el promedio simple de $DRES$ durante
	1990--1995 sea igual o superior al 20\,\%; por tanto, la pertenencia a la
	muestra se determina antes del período econométrico 1996--2021.

	\renewcommand{\arraystretch}{1.05}
	\setlength{\tabcolsep}{5pt}

	\begin{footnotesize}
		\begin{longtable}{
				L{5.1cm}
				>{\RaggedLeft\arraybackslash}p{1.3cm}
				L{5.1cm}
				>{\RaggedLeft\arraybackslash}p{1.3cm}
			}

			\caption[Países seleccionados mediante DRES]{Países seleccionados mediante el umbral de $DRES\geq 20\%$.}
			\label{tab:muestra_dres20} \\
			\hline
			País (\isocode{ISO3}) & DRES medio & País (\isocode{ISO3}) & DRES medio \\
			\hline
			\endfirsthead

			\hline
			País (\isocode{ISO3}) & DRES medio & País (\isocode{ISO3}) & DRES medio \\
			\hline
			\endhead

			Angola (\isocode{AGO}) & 98,97\,\% & Albania (\isocode{ALB}) & 21,41\,\% \\
			Emiratos Árabes Unidos (\isocode{ARE}) & 85,44\,\% & Antigua y Barbuda (\isocode{ATG}) & 29,09\,\% \\
			Australia (\isocode{AUS}) & 46,44\,\% & Burundi (\isocode{BDI}) & 32,85\,\% \\
			Baréin (\isocode{BHR}) & 71,82\,\% & Bolivia (\isocode{BOL}) & 56,92\,\% \\
			Brunéi (\isocode{BRN}) & 86,42\,\% & República Centroafricana (\isocode{CAF}) & 61,56\,\% \\
			Chile (\isocode{CHL}) & 48,37\,\% & Camerún (\isocode{CMR}) & 47,60\,\% \\
			República Democrática del Congo (\isocode{COD}) & 81,42\,\% & República del Congo (\isocode{COG}) & 84,66\,\% \\
			Colombia (\isocode{COL}) & 30,99\,\% & Argelia (\isocode{DZA}) & 82,02\,\% \\
			Ecuador (\isocode{ECU}) & 38,21\,\% & Egipto (\isocode{EGY}) & 38,34\,\% \\
			Gabón (\isocode{GAB}) & 84,60\,\% & Ghana (\isocode{GHA}) & 43,98\,\% \\
			Guinea (\isocode{GIN}) & 88,57\,\% & Gambia (\isocode{GMB}) & 55,64\,\% \\
			Guyana (\isocode{GUY}) & 43,59\,\% & Indonesia (\isocode{IDN}) & 37,21\,\% \\
			India (\isocode{IND}) & 21,51\,\% & Irán (\isocode{IRN}) & 88,31\,\% \\
			Irak (\isocode{IRQ}) & 97,58\,\% & Israel (\isocode{ISR}) & 33,23\,\% \\
			Jamaica (\isocode{JAM}) & 44,88\,\% & Jordania (\isocode{JOR}) & 25,40\,\% \\
			Kuwait (\isocode{KWT}) & 89,22\,\% & Liberia (\isocode{LBR}) & 49,88\,\% \\
			Libia (\isocode{LBY}) & 92,80\,\% & Mongolia (\isocode{MNG}) & 33,79\,\% \\
			Mauritania (\isocode{MRT}) & 52,55\,\% & Nigeria (\isocode{NGA}) & 93,61\,\% \\
			Noruega (\isocode{NOR}) & 56,08\,\% & Nauru (\isocode{NRU}) & 88,39\,\% \\
			Omán (\isocode{OMN}) & 82,80\,\% & Perú (\isocode{PER}) & 51,55\,\% \\
			Papúa Nueva Guinea (\isocode{PNG}) & 63,63\,\% & Catar (\isocode{QAT}) & 82,28\,\% \\
			Rusia (\isocode{RUS}) & 54,91\,\% & Arabia Saudita (\isocode{SAU}) & 87,39\,\% \\
			Sierra Leona (\isocode{SLE}) & 76,97\,\% & Surinam (\isocode{SUR}) & 75,03\,\% \\
			Seychelles (\isocode{SYC}) & 22,82\,\% & Siria (\isocode{SYR}) & 59,00\,\% \\
			Togo (\isocode{TGO}) & 43,63\,\% & Trinidad y Tobago (\isocode{TTO}) & 56,20\,\% \\
			Venezuela (\isocode{VEN}) & 82,95\,\% & Vietnam (\isocode{VNM}) & 28,76\,\% \\
			Yemen (\isocode{YEM}) & 93,52\,\% & Sudáfrica (\isocode{ZAF}) & 41,72\,\% \\
			Zambia (\isocode{ZMB}) & 91,12\,\% & & \\

			\hline
			\multicolumn{4}{p{14.1cm}}{\textit{Nota: Los países se ordenan por código \isocode{ISO3} de izquierda a derecha y de arriba abajo. \isocode{COD} y \isocode{COG} representan dos países distintos: República Democrática del Congo y República del Congo, respectivamente. $DRES$ medio corresponde al promedio simple de las seis observaciones anuales de 1990--1995. Todos los países satisfacen $DRES\geq 20\%$ y cuentan con información comercial completa durante el período base.}} \\
		\end{longtable}
	\end{footnotesize}
\thesissource{Elaboración propia a partir de \texttt{data/processed/dres/dres\_sample\_20.csv}.}

	\section{Figuras conceptuales y esquemas analíticos}
	\label{sec:figuras}

	Esta sección reúne figuras conceptuales y esquemas analíticos que complementan la exposición teórica y metodológica del trabajo. Los diagramas permiten sintetizar los principales mecanismos de transición estructural y las relaciones entre las variables centrales del estudio, facilitando la lectura del enfoque general de la tesis.

	\begin{figure}[htbp]
		\centering
		\resizebox{0.95\textwidth}{!}{%
			\begin{tikzpicture}[
				font=\small,
				box/.style={
					draw,
					rectangle,
					rounded corners=6pt,
					align=center,
					minimum width=12.4cm,
					inner sep=6pt,
					fill=gray!4,
					line width=0.85pt
				},
				sidebox/.style={
					draw,
					rectangle,
					rounded corners=6pt,
					align=center,
					minimum width=5.4cm,
					inner sep=6pt,
					fill=gray!2,
					line width=0.85pt
				},
				arrow/.style={-{Stealth[length=3mm]}, line width=0.9pt},
				dasharrow/.style={-{Stealth[length=3mm]}, line width=0.9pt, dashed}
				]

				% Nivel 1: Selección (y = 7.5)
				\node[box] (dres) at (0, 7.5) {\textbf{Criterio de Selección Muestral}\\[3pt]
					$DRES \ge 20\%$: 55 economías iniciales; muestra común: 49 economías (1996--2021)};

				% Nivel 2: Mecanismo focal (y = 5.0)
				\node[box] (rents) at (0, 5.0) {\textbf{Exposición focal y heterogeneidad sectorial}\\[3pt]
					Rentas Extractivas Totales ($RENTS$)\\[2pt]
					\footnotesize Hidrocarburos ($rents\_oil\_gas$) \quad $\cdot$ \quad Minería y Carbón ($rents\_mining$)};

				% Nivel 3: Canales explicativos (y = 2.4, x=-3.6 y x=+3.6 -> deja 1.8cm libres al centro)
				\node[sidebox] (inst) at (-3.6, 2.4) {\textbf{Canal Institucional}\\[3pt]
					$ECI$: $INST$ ($\beta = +0,267^*$)\\[2pt]
					\footnotesize Interacción $RENTS \times INST$: No signif.};

				\node[sidebox] (cap) at (3.6, 2.4) {\textbf{Capacidades y Estructura}\\[3pt]
					$DIVX$: $HUMCAP$ ($\beta = +0,107^{**}$)\\[2pt]
					\footnotesize $ECI$: $FEXP$ ($\beta = -0,0068^{***}$)};

				% Nivel 4: Resultados (y = -0.3)
				\node[box] (out) at (0, -0.3) {\textbf{Resultados de Transformación Estructural Externa}\\[4pt]
					Complejidad cualitativa ($ECI, \beta_{RENTS} = -0,010^{**}, p = 0,023$)\\[2pt]
					Diversificación cuantitativa ($DIVX, \beta_{RENTS} = -0,004^{***}, p = 0,009$)};

				% Flechas
				\draw[arrow] (dres) -- (rents);
				\draw[arrow] (rents) -- (out);

				\draw[dasharrow] (inst) -- (out);
				\draw[arrow] (cap) -- (out);

			\end{tikzpicture}%
		}

		\caption[Diseño empírico y asociaciones principales]{Diseño empírico y asociaciones principales entre la selección por dependencia extractiva ($DRES$), las rentas ($RENTS$), los canales considerados y los resultados de complejidad ($ECI$) y diversificación ($DIVX$).}

		\label{fig:arquitectura-mecanismos}
		\caption*{\footnotesize \textit{Nota:} Las flechas organizan el diseño analítico y no representan relaciones causales. \textit{Fuente: Elaboración propia a partir de las estimaciones TWFE presentadas en el Capítulo 11.}}
	\end{figure}

	Con el fin de evaluar un resultado complementario de transformación estructural externa, la tesis contrasta el modelo basado en el \textit{Economic Complexity Index} (ECI) con una segunda especificación en la que la variable dependiente corresponde a la diversificación exportadora ($DIVX=1-HHI$). Esta medida no reemplaza al ECI, sino que permite observar si los mecanismos asociados a la sofisticación de la canasta exportadora también se relacionan con una menor concentración de dicha canasta. En el modelo con $DIVX$, $HHI$ se excluye del conjunto de regresores para evitar una relación mecánica entre la variable dependiente y sus componentes de construcción.



	\section{Gráficos complementarios}
	\label{sec:plots}

	\subsection{Perfiles extractivos e ingreso en el panel empírico}
	\label{sec:figuras-panel-ingreso}

	Como análisis descriptivo de la heterogeneidad de las economías seleccionadas en el panel ($DRES \ge 20\%$), la \autoref{fig:panel-commodities-income} presenta la distribución de la dependencia exportadora extractiva ($DRES\%$) clasificada según el grupo de ingreso del Banco Mundial: bajos ingresos (\textit{LIC}), ingresos medios bajos (\textit{LMIC}), ingresos medios altos (\textit{UMIC}) y altos ingresos (\textit{HIC}).

	La \autoref{fig:panel-commodities-income} muestra que la dependencia exportadora extractiva no disminuye de manera monotónica con el nivel de ingreso per cápita. Los grupos de ingreso medio alto (\textit{UMIC}) y alto (\textit{HIC}) también albergan economías con una elevada especialización en bienes primarios no renovables del subsuelo (tales como Noruega, Australia, Kuwait, Omán o Qatar).

	\begin{figure}[htbp]
		\centering
		\includegraphics[width=0.85\textwidth]{\subfix{../../figures/panel_extractive_dependence_by_income_group}}
		\caption[Dependencia exportadora extractiva por grupo de ingreso en el panel]{Distribución de la dependencia exportadora extractiva ($DRES\%$) en las economías del panel empírico según grupo de ingreso: bajos ingresos (LIC), ingresos medios bajos (LMIC), ingresos medios altos (UMIC) y altos ingresos (HIC).}
		\label{fig:panel-commodities-income}
		\caption*{\footnotesize \textit{Fuente: Elaboración propia a partir de $DRES$ medio durante 1990--1995 y la clasificación de ingresos utilizada en la investigación.}}
	\end{figure}

	De manera específica para el sector de hidrocarburos, la \autoref{fig:panel-oil-gas-rents-income} analiza la distribución de las rentas de petróleo y gas natural ($RENTS^{OIL+GAS}$ como $\%$ del PIB) a través de los cuatro grupos de ingreso. Se observan valores elevados en varias economías de ingresos altos (\textit{HIC}) y medios altos (\textit{UMIC}); esta comparación es descriptiva y no identifica el efecto de las rentas sobre el nivel de ingreso.

	\begin{figure}[htbp]
		\centering
		\includegraphics[width=0.85\textwidth]{\subfix{../../figures/panel_oil_gas_rents_by_income_group}}
		\caption[Rentas de petróleo y gas natural por grupo de ingreso en el panel]{Distribución de las rentas de petróleo y gas natural ($RENTS^{OIL+GAS}$ como $\%$ del PIB) en las economías del panel empírico según grupo de ingreso: bajos ingresos (LIC), ingresos medios bajos (LMIC), ingresos medios altos (UMIC) y altos ingresos (HIC).}
		\label{fig:panel-oil-gas-rents-income}
		\caption*{\footnotesize \textit{Fuente: Elaboración propia a partir del panel maestro de la investigación (1996--2021).}}
	\end{figure}

	Por su parte, la \autoref{fig:panel-mining-rents-income} presenta la distribución de las rentas mineras ($RENTS^{MINING}$ como $\%$ del PIB) a través de los cuatro grupos de ingreso. En la muestra, los valores más elevados se observan principalmente en economías de ingresos medios bajos (\textit{LMIC}) y medios altos (\textit{UMIC}); el gráfico no implica que el grupo de ingreso determine la composición extractiva.

	\begin{figure}[htbp]
		\centering
		\includegraphics[width=0.85\textwidth]{\subfix{../../figures/panel_mining_rents_by_income_group}}
		\caption[Rentas de minería por grupo de ingreso en el panel]{Distribución de las rentas de minería ($RENTS^{MINING}$ como $\%$ del PIB) en las economías del panel empírico según grupo de ingreso: bajos ingresos (LIC), ingresos medios bajos (LMIC), ingresos medios altos (UMIC) y altos ingresos (HIC).}
		\label{fig:panel-mining-rents-income}
		\caption*{\footnotesize \textit{Fuente: Elaboración propia a partir del panel maestro de la investigación (1996--2021).}}
	\end{figure}

	Como análisis complementario de la heterogeneidad de las economías extractivas, la \autoref{fig:recreated_owjimehr} presenta la evolución temporal del Índice de Complejidad Económica (ECI) para una selección de los principales países exportadores de recursos de la muestra (1996--2021) (réplica metodológica adaptada de \textcite{Owjimehr2024}). A diferencia del gráfico original de rankings, en este caso se presenta la trayectoria del puntaje de ECI continuo, lo que permite apreciar las brechas absolutas de capacidades productivas de manera más intuitiva. Se observa una notable persistencia en los extremos: Noruega mantiene de forma constante los mayores niveles de complejidad, mientras que economías como Angola y el Congo permanecen en niveles sumamente bajos durante todo el período.

	\begin{figure}[htbp]
		\centering
		\includegraphics[width=0.85\textwidth]{\subfix{../../figures/recreated_owjimehr}}
		\caption[Evolución del ECI en países exportadores de recursos]{Evolución del Índice de Complejidad Económica (ECI) en una selección de los principales países exportadores de recursos, 1996--2021 (adaptado de \textcite{Owjimehr2024}).}
		\label{fig:recreated_owjimehr}
		\caption*{\footnotesize \textit{Fuente: Elaboración propia a partir del panel maestro de la investigación, adaptando el análisis de \textcite{Owjimehr2024}.}}
	\end{figure}

	\section{Mapas}
	\label{sec:mapas}

	La \autoref{fig:resource-rents-world-map} muestra la distribución global de las rentas de recursos naturales como porcentaje del PIB.

	\begin{figure}[htbp]
		\centering
		\includegraphics[width=\textwidth]{\subfix{../../figures/resource_rents_world_map}}

		\caption[Rentas de recursos naturales (\% del PIB)]{Distribución global de las rentas de recursos naturales como porcentaje del PIB.}

		\label{fig:resource-rents-world-map}
		\caption*{\footnotesize \textit{Fuente: Elaboración propia con base en el panel de datos de la investigación.}}
	\end{figure}

	La \autoref{fig:eci-world-map} presenta la distribución global del Índice de Complejidad Económica (ECI), un indicador que captura el grado de sofisticación productiva de las economías a partir de la diversidad y ubicuidad de los bienes que exportan. En general, valores más altos del ECI se asocian con estructuras productivas más diversificadas y mayores capacidades tecnológicas, mientras que valores bajos reflejan economías más concentradas en productos primarios o de menor intensidad tecnológica; en este contexto, el mapa ofrece un marco descriptivo para observar las diferencias en capacidades productivas entre países y su posible relación con la dependencia de recursos naturales.

	\begin{figure}[htbp]
		\centering
		\includegraphics[width=\textwidth]{\subfix{../../figures/eci_world_map}}

		\caption[Distribución global del ECI]{Distribución global del Índice de Complejidad Económica (ECI).}

		\label{fig:eci-world-map}
		\caption*{\footnotesize \textit{Fuente: Elaboración propia con base en el panel de datos de la investigación.}}
	\end{figure}

	La \autoref{fig:eci-world-map-quartiles} muestra la distribución global del Índice de Complejidad Económica (ECI) clasificada por cuartiles. Esta representación permite comparar de manera relativa el nivel de sofisticación productiva entre economías, agrupando a los países según su posición en la distribución mundial del indicador. En particular, el primer cuartil (Q1) corresponde a economías con menor complejidad económica, mientras que el cuarto cuartil (Q4) agrupa a las economías con estructuras productivas más diversificadas y tecnológicamente avanzadas.

	\begin{figure}[htbp]
		\centering
		\includegraphics[width=\textwidth]{\subfix{../../figures/eci_world_map_quartiles}}

		\caption[ECI por cuartiles]{Distribución global del Índice de Complejidad Económica (ECI) por cuartiles.}

		\label{fig:eci-world-map-quartiles}
		\caption*{\footnotesize \textit{Fuente: Elaboración propia con base en el panel de datos de la investigación.}}
	\end{figure}


	Una representación especialmente ilustrativa de la estructura productiva mundial es el denominado \textit{Product Space}, desarrollado por \textcite{Hausmann2014}. Esta red vincula productos según su proximidad productiva, definida a partir de la probabilidad de que dos bienes sean exportados conjuntamente por los mismos países. La visualización resultante revela una topología altamente estructurada en la que los bienes primarios y los \textit{commodities} tienden a ubicarse en la periferia, mientras que las manufacturas complejas —como maquinaria, productos químicos y electrónica— ocupan el núcleo densamente interconectado. En este marco, el desarrollo económico puede entenderse como un proceso de acumulación gradual de capacidades productivas que permite a los países desplazarse desde actividades periféricas hacia sectores más centrales y complejos.

	\begin{figure}[htbp]
		\centering
		\includegraphics[width=0.95\textwidth]{\subfix{../../figures/3_hausmann_2014.png}}

		\caption[El Product Space]{El \textit{Product Space}: red global de proximidad entre productos basada en patrones de coexportación entre países. Los nodos representan productos y los enlaces reflejan proximidad productiva.}

		\label{fig:product_space_hausmann}
		\caption*{\footnotesize \textit{Fuente: Adaptado de \textcite{Hausmann2014}.}}
	\end{figure}

	La literatura sobre la \textit{Dutch disease} ha señalado que los auges en los sectores de recursos naturales pueden generar apreciaciones del tipo de cambio real, reasignación de factores productivos y pérdida de competitividad en sectores transables, particularmente en la manufactura \parencite{Barder2006}. Desde la perspectiva del \textit{Product Space} propuesto por \textcite{Hausmann2014}, este fenómeno puede interpretarse como un desplazamiento estructural de la economía hacia regiones periféricas del espacio productivo. En dicho espacio, las actividades de manufactura —especialmente aquellas asociadas a productos más sofisticados— tienden a ubicarse en el núcleo densamente interconectado de la red, debido a la gran cantidad de capacidades productivas compartidas que requieren. En contraste, los bienes primarios y extractivos suelen localizarse en la periferia, caracterizada por un menor número de conexiones con otras actividades productivas. En este sentido, una especialización excesiva en recursos naturales no solo debilita la base manufacturera de la economía, sino que también reduce las oportunidades cercanas para diversificación hacia productos más complejos. En consecuencia, la \textit{Dutch disease} no solo tiene implicaciones macroeconómicas de corto plazo, sino que también puede afectar la trayectoria de acumulación de capacidades productivas, limitando la capacidad de la economía para transitar hacia sectores de manufactura ubicados en el núcleo del espacio productivo.

	Como complemento visual al marco analítico de la tesis, la \autoref{fig:product_space_communities} muestra cómo distintos grupos de productos se ubican en posiciones heterogéneas en términos de complejidad económica y conectividad productiva. En particular, los sectores asociados a manufacturas y bienes industriales tienden a concentrarse en regiones más densamente conectadas del espacio productivo, mientras que muchas actividades primarias y extractivas aparecen en posiciones más periféricas. Esta configuración sugiere que las economías especializadas en recursos naturales pueden enfrentar mayores restricciones estructurales para diversificar hacia actividades de mayor complejidad.

	\begin{figure}[htbp]
		\centering
		\includegraphics[width=0.9\textwidth]{\subfix{../../figures/4_hausmann_2014.png}}

		\caption[Complejidad y conectividad en el Product Space]{Complejidad promedio y conectividad de comunidades de productos en el espacio productivo global. El tamaño de las burbujas es proporcional a la participación de cada comunidad en el comercio mundial.}

		\label{fig:product_space_communities}
		\caption*{\footnotesize \textit{Fuente: Adaptado de \textcite{Hausmann2014}, The Atlas of Economic Complexity.}}
	\end{figure}

	Como se observa en la \autoref{tab:product_communities_selected}, las comunidades asociadas a industrias extractivas presentan valores sustancialmente menores del \textit{Product Complexity Index} (PCI) en comparación con sectores manufactureros intensivos en conocimiento. Este contraste refleja diferencias estructurales en las capacidades productivas requeridas por cada tipo de actividad y resulta consistente con la literatura sobre complejidad económica y diversificación productiva \parencite{Hausmann2014}.

	\begin{table}[htbp]
		\centering
		\small
		\caption[Comunidades de productos y complejidad (PCI)]{
			Características seleccionadas de comunidades de productos en el espacio productivo, incluyendo complejidad promedio (PCI), número de productos y participación en el comercio mundial.
		}
		\label{tab:product_communities_selected}

		\begin{threeparttable}
			\begin{adjustbox}{max width=\textwidth}
			\begin{tabular}{L{2.8cm}L{2.2cm}L{2.6cm}L{2.2cm}L{2.8cm}}
				\toprule
				Comunidad & PCI promedio & Número de productos & Comercio mundial & Participación mundial \\
				\midrule
				\multicolumn{5}{l}{\textit{Industrias extractivas}} \\

				Petróleo & -2.08 & 4 & 2.3T & 10.49\% \\
				Minería & -0.59 & 48 & 1.1T & 5.01\% \\
				Carbón & 0.21 & 6 & 183B & 0.83\% \\

				\midrule
				\multicolumn{5}{l}{\textit{Sectores manufactureros (contraste)}} \\

				Maquinaria & 2.54 & 125 & 4.4T & 20.29\% \\
				Electrónica & 2.25 & 52 & 3.6T & 16.71\% \\
				Químicos \& salud & 2.52 & 64 & 1.6T & 7.47\% \\

				\bottomrule
			\end{tabular}
			\end{adjustbox}

			\begin{tablenotes}
				\footnotesize
				\item[] \centering\textit{Fuente: Elaboración propia a partir de datos reportados en \textcite{Hausmann2014}, The Atlas of Economic Complexity.}
			\end{tablenotes}

		\end{threeparttable}
	\end{table}


	Las visualizaciones del espacio productivo pueden consultarse de forma interactiva en el \textit{Atlas of Economic Complexity}, desarrollado por el Harvard Growth Lab \parencite{HarvardGrowthLab2024}.



	\section{Muestra analítica y diagnósticos del panel}
	\label{sec:anexo_diagnosticos_series_tiempo}

	Los diagnósticos de este anexo corresponden a la muestra econométrica común
	y congelada de la investigación. No se utilizan para seleccionar automáticamente un estimador ni
	para clasificar todas las series como $I(0)$ o $I(1)$; su función es
	documentar la estructura temporal y residual de las especificaciones TWFE.
	La \autoref{tab:anexo_contrato_muestra_peer2} resume el tránsito desde la
	grilla inicial hasta la muestra común de casos completos utilizada en los
	modelos de $ECI$ y $DIVX$.

	\begin{table}[htbp]
		\centering
		\small
		\caption{Contrato de la muestra econométrica común}
		\label{tab:anexo_contrato_muestra_peer2}
		\begin{tabular}{lrr}
			\toprule
			Dimensión & Grilla inicial & Muestra econométrica \\
			\midrule
			Observaciones país--año & 1.430 & 1.044 \\
			Países & 55 & 49 \\
			Años calendario, 1996--2021 & 26 & 26 \\
			Años efectivos con observaciones & 26 & 23 \\
			Años con cobertura completa & --- & 6 \\
			Años con cobertura parcial & --- & 17 \\
			Años sin observaciones comunes & --- & 3 \\
			\bottomrule
		\end{tabular}
		\caption*{\footnotesize \textit{Nota:} Los años sin observaciones comunes son 1997, 1999 y 2001. La muestra es idéntica para $ECI$ y $DIVX$, sigue la regla de casos completos y no utiliza imputación ni interpolación.}
		\caption*{\footnotesize \textit{Fuente: Elaboración propia a partir de la validación de la muestra econométrica.}}
	\end{table}

	La persistencia AR(1) de la \autoref{tab:anexo_persistencia_fisher_peer2}
	se estima de manera descriptiva sobre 869 pares temporales válidos. La misma
	tabla resume cuántas de ocho variantes Fisher--ADF rechazan al 5\% la
	hipótesis nula de que todos los paneles contienen raíz unitaria.

	\begin{table}[htbp]
		\centering
		\small
		\caption{Persistencia temporal y síntesis de pruebas Fisher--ADF}
		\label{tab:anexo_persistencia_fisher_peer2}
		\begin{tabular}{lrrr}
			\toprule
			Variable & AR(1) & Rechazos en niveles & Rechazos en diferencias \\
			\midrule
			$ECI$   & 0,495 & 6 de 8 & 8 de 8 \\
			$DIVX$  & 0,744 & 1 de 8 & 8 de 8 \\
			$RENTS$ & 0,701 & 1 de 8 & 6 de 8 \\
			$INST$  & 0,809 & 6 de 8 & 8 de 8 \\
			\bottomrule
		\end{tabular}
		\caption*{\footnotesize \textit{Nota:} Los coeficientes AR(1) son descriptivos y se estiman condicionando por efectos fijos de país y año, con errores estándar agrupados por país. Las variantes Fisher--ADF combinan constante o tendencia, uno o dos rezagos y series originales o descentradas transversalmente. Fisher--ADF permite los huecos internos de este panel; IPS no resulta aplicable bajo esa estructura. Un rechazo solo indica que al menos uno de los paneles es estacionario y no permite asignar un orden de integración común.}
		\caption*{\footnotesize \textit{Fuente: Elaboración propia a partir de los diagnósticos temporales de la investigación.}}
	\end{table}

	La \autoref{tab:anexo_diagnosticos_m3_peer2} reúne los contrastes aplicados a
	los residuos y regresores de $M_3$. Los resultados rechazan homocedasticidad y
	ausencia de autocorrelación en ambas ecuaciones, no rechazan independencia
	transversal y muestran niveles bajos de colinealidad entre los regresores.

	\begin{table}[htbp]
		\centering
		\small
		\caption{Diagnósticos residuales y de colinealidad de $M_3$}
		\label{tab:anexo_diagnosticos_m3_peer2}
		\begin{adjustbox}{max width=\textwidth}
		\begin{tabular}{llrrl}
			\toprule
			Modelo & Diagnóstico & Estadístico & $p$-valor & Lectura \\
			\midrule
			$ECI$ & Wald modificado & 9.046,05 & $<0,001$ & Rechaza homocedasticidad \\
			$ECI$ & Wooldridge AR(1) & 23,95 & $<0,001$ & Rechaza ausencia de autocorrelación \\
			$ECI$ & Pesaran CD & $-0,44$ & 0,663 & No rechaza independencia transversal \\
			$ECI$ & VIF \textit{within} máximo & 1,956 & --- & Bajo \\
			\addlinespace
			$DIVX$ & Wald modificado & 13.046,62 & $<0,001$ & Rechaza homocedasticidad \\
			$DIVX$ & Wooldridge AR(1) & 82,68 & $<0,001$ & Rechaza ausencia de autocorrelación \\
			$DIVX$ & Pesaran CD & $-0,77$ & 0,442 & No rechaza independencia transversal \\
			$DIVX$ & VIF \textit{within} máximo & 1,929 & --- & Bajo \\
			\bottomrule
		\end{tabular}
		\end{adjustbox}
		\caption*{\footnotesize \textit{Nota:} Los contrastes utilizan 1.044 observaciones y 49 países. Los VIF se calculan después de retirar los efectos fijos de país y año.}
		\caption*{\footnotesize \textit{Fuente: Elaboración propia a partir de los diagnósticos del modelo completo.}}
	\end{table}

	\section{Sensibilidades de las estimaciones principales}
	\label{sec:anexo_sensibilidades_peer2}

	La \autoref{tab:anexo_robustez_focal_peer2} contrasta los términos focales de
	$M_3$ con bootstrap \textit{wild cluster}, exclusión iterativa de países y
	retiro de observaciones influyentes. Estos ejercicios complementan la
	inferencia agrupada por país; no sustituyen el modelo principal.

	\begin{table}[htbp]
		\centering
		\footnotesize
		\caption{Robustez de los términos focales de $M_3$}
		\label{tab:anexo_robustez_focal_peer2}
		\begin{adjustbox}{max width=\textwidth}
		\begin{tabular}{llrrrrrr}
			\toprule
			Modelo & Término & $\hat\beta$ & $p$ & $p_{boot}$ & Cambios de signo LOO & Sig. LOO & $\hat\beta$ sin alertas ($p$) \\
			\midrule
			$ECI$ & $RENTS$ & $-0,0098$ & 0,023 & 0,040 & 0 & 46/49 al 5\% & $-0,0053$ (0,046) \\
			$ECI$ & $RENTS\times INST$ & $-0,0052$ & 0,139 & 0,209 & 0 & 4/49 al 10\% & $-0,0037$ (0,147) \\
			$DIVX$ & $RENTS$ & $-0,0035$ & 0,009 & 0,021 & 0 & 49/49 al 5\% & $-0,0045$ ($<0,001$) \\
			$DIVX$ & $RENTS\times INST$ & 0,0009 & 0,427 & 0,477 & 0 & 0/49 al 10\% & 0,0004 (0,694) \\
			\bottomrule
		\end{tabular}
		\end{adjustbox}
		\caption*{\footnotesize \textit{Nota:} El bootstrap emplea 9.999 replicaciones y pesos Rademacher. LOO indica la exclusión de un país por vez. La columna final excluye las observaciones identificadas por los diagnósticos de influencia.}
		\caption*{\footnotesize \textit{Fuente: Elaboración propia a partir de los ejercicios de estabilidad de las estimaciones.}}
	\end{table}

	La \autoref{tab:anexo_tendencias_2014_peer2} muestra que las tendencias
	lineales por país conservan el signo de $RENTS$ en ambas
	ecuaciones, aunque su precisión se mantiene en $DIVX$ y se pierde en $ECI$.
	La sensibilidad temporal de 2014 se estima sobre la muestra completa mediante
	interacciones, no mediante dos regresiones por submuestras. Por ello se
	interpreta como un contraste de estabilidad temporal y no como una prueba de
	quiebre estructural.

	\begin{table}[htbp]
		\centering
		\small
		\caption{Tendencias por país y sensibilidad temporal de 2014}
		\label{tab:anexo_tendencias_2014_peer2}
		\begin{adjustbox}{max width=\textwidth}
		\begin{tabular}{llrrl}
			\toprule
			Modelo & Contraste & Coeficiente & $p$-valor & Lectura \\
			\midrule
			$ECI$ & $RENTS$, con tendencias por país & $-0,0044$ & 0,210 & Signo estable; menor precisión \\
			$ECI$ & $RENTS\times INST$, con tendencias & 0,0006 & 0,822 & Sin evidencia de moderación \\
			$ECI$ & Cambio conjunto desde 2014 & --- & 0,343 & No rechaza estabilidad focal \\
			\addlinespace
			$DIVX$ & $RENTS$, con tendencias por país & $-0,0030$ & 0,020 & Signo y precisión estables \\
			$DIVX$ & $RENTS\times INST$, con tendencias & $-0,0003$ & 0,855 & Sin evidencia de moderación \\
			$DIVX$ & Cambio conjunto desde 2014 & --- & 0,059 & Sensibilidad marginal al 10\% \\
			\bottomrule
		\end{tabular}
		\end{adjustbox}
		\caption*{\footnotesize \textit{Nota:} El contraste de 2014 prueba conjuntamente los cambios en $RENTS$, $INST$ y $RENTS\times INST$ dentro de una ecuación única con 1.044 observaciones. Los datos anuales de 2014 combinan meses anteriores y posteriores a la caída del precio del petróleo de mitad de año.}
		\caption*{\footnotesize \textit{Fuente: Elaboración propia a partir de los ejercicios de sensibilidad temporal.}}
	\end{table}

	La \autoref{tab:anexo_efectos_marginales_peer2} presenta el efecto marginal
	de $RENTS$ en tres percentiles representativos de $INST$. Aunque la pendiente
	es negativa en los seis puntos reportados, su precisión y su variación a lo
	largo del soporte institucional difieren entre $ECI$ y $DIVX$; por ello, estos
	resultados no constituyen evidencia de una moderación institucional general.

	\begin{table}[htbp]
		\centering
		\small
		\caption{Efecto marginal de $RENTS$ dentro del soporte de $INST$}
		\label{tab:anexo_efectos_marginales_peer2}
		\begin{tabular}{llrr}
			\toprule
			Modelo & Punto de $INST$ & Efecto marginal & $p$-valor \\
			\midrule
			$ECI$ & Percentil 10 & $-0,0034$ & 0,448 \\
			$ECI$ & Percentil 50 & $-0,0078$ & 0,039 \\
			$ECI$ & Percentil 90 & $-0,0146$ & 0,021 \\
			\addlinespace
			$DIVX$ & Percentil 10 & $-0,0047$ & $<0,001$ \\
			$DIVX$ & Percentil 50 & $-0,0039$ & $<0,001$ \\
			$DIVX$ & Percentil 90 & $-0,0027$ & 0,214 \\
			\bottomrule
		\end{tabular}
		\caption*{\footnotesize \textit{Nota:} Los efectos se calculan en cinco puntos entre los percentiles 10 y 90; se reportan tres puntos representativos. Las pendientes no son significativas en todo el soporte y siguen patrones opuestos entre ecuaciones, por lo que no respaldan una moderación institucional general.}
		\caption*{\footnotesize \textit{Fuente: Elaboración propia a partir de las estimaciones de efectos marginales.}}
	\end{table}

	\section{Desagregación complementaria de las rentas extractivas}
	\label{sec:anexo_desagregacion_peer2}

	La desagregación sustituye $RENTS$ por rentas de petróleo y gas y por rentas
	de minería y carbón, manteniendo los controles, los efectos fijos y la muestra
	de $M_3$. La \autoref{tab:anexo_desagregacion_peer2} resume los coeficientes
	directos y su estabilidad. Las cuatro interacciones con $INST$ carecen de
	precisión convencional y con bootstrap.

	\begin{table}[htbp]
		\centering
		\small
		\caption{Resultados complementarios por componente extractivo}
		\label{tab:anexo_desagregacion_peer2}
		\begin{adjustbox}{max width=\textwidth}
		\begin{tabular}{llrrrrl}
			\toprule
			Modelo & Componente & $\hat\beta$ & $p$ & $p_{boot}$ & Cambios de signo LOO & Clasificación \\
			\midrule
			$ECI$ & Petróleo y gas & $-0,0073$ & 0,212 & 0,292 & 0 & No concluyente \\
			$ECI$ & Minería y carbón & $-0,0136$ & 0,001 & 0,024 & 0 & Complementaria \\
			$DIVX$ & Petróleo y gas & $-0,0048$ & 0,001 & 0,003 & 0 & Complementaria \\
			$DIVX$ & Minería y carbón & $-0,0014$ & 0,459 & 0,542 & 1 & No concluyente \\
			\bottomrule
		\end{tabular}
		\end{adjustbox}
		\caption*{\footnotesize \textit{Nota:} Los dos resultados clasificados como complementarios conservan el signo en las 49 exclusiones de países y mantienen significación con bootstrap. La desagregación no reemplaza el resultado agregado principal.}
		\caption*{\footnotesize \textit{Fuente: Elaboración propia a partir de la estimación desagregada de las rentas extractivas.}}
	\end{table}

	\thesisendchap{anexos}

\end{document}
